\documentclass[12pt]{article}
\def\activityid{F07-M-v2}\usepackage[letterpaper,margin=.65in,top=.60in,bottom=.65in]{geometry}
\usepackage[T1]{fontenc}\usepackage[default]{sourcesanspro}
\usepackage{microtype,tikz,array,tabularx,fancyhdr,amsmath}\usepackage[hidelinks]{hyperref}
\usetikzlibrary{arrows.meta,calc}
\setlength{\parindent}{0pt}\setlength{\parskip}{7pt}\setlength{\headheight}{0pt}\setlength{\footskip}{26pt}\setlength{\emergencystretch}{2em}
\pagestyle{fancy}\fancyhf{}\renewcommand{\headrulewidth}{0pt}\renewcommand{\footrulewidth}{.4pt}
\fancyfoot[L]{\scriptsize Bellingham Math Circle / Week 7 / \activityid}\fancyfoot[R]{\scriptsize \thepage}
\definecolor{ink}{gray}{.12}\definecolor{soft}{gray}{.95}\color{ink}
\newcommand{\PageTitle}[2]{{\fontsize{10}{12}\selectfont\bfseries WEEK 7\quad GAME LAB\hfill #1\par}\vspace{3pt}{\fontsize{26}{29}\selectfont\bfseries #2\par}\vspace{2pt}}
\newcommand{\NameLine}{{\small Name: \rule{2.65in}{.4pt}\hfill Date: \rule{1.35in}{.4pt}\par}\vspace{4pt}}
\newcommand{\Task}[2]{\par\vspace{3pt}{\large\bfseries #1}\par #2\par}
\newcommand{\Subhead}[1]{\par\vspace{3pt}{\large\bfseries #1}\par}
\newcommand{\RuleBox}[1]{\par\begingroup\setlength{\fboxsep}{8pt}\colorbox{soft}{\parbox{\dimexpr\linewidth-16pt\relax}{#1}}\endgroup\par}
\newcommand{\WriteBox}[2]{\par\textbf{#1}\par\vspace{2pt}\begin{tikzpicture}\draw[black!40,line width=.5pt] (0,0) rectangle (\dimexpr\linewidth-.5pt\relax,#2);\end{tikzpicture}\par}
\newcommand{\StudentFooter}{\vfill{\small Build first. Explain by pointing, talking, or drawing.}\par}
\newcommand{\Code}[1]{\texttt{#1}}

\hypersetup{pdftitle={Week 7: grades-2-3}}
\begin{document}
\PageTitle{GRADES 2--3}{The missing move changes everything}\NameLine
\RuleBox{Two players alternate. Take \textbf{1, 3, or 4} counters each turn. Taking 2 is illegal. Never take more than remain. Taking the last counter wins.}\Task{1. Play and suspect.}{Try piles of 5, 6, and 7. Trade who goes first. A win gives you a clue; it does not yet prove a strategy.}
\Task{2. Solve small positions.}{L means the player about to move loses against perfect replies. W means that player has a winning strategy. Zero is L: the previous player took the last counter.}
\begin{center}\renewcommand{\arraystretch}{1.85}\begin{tabular}{p{.68in}|p{.68in}|p{.68in}|p{.68in}|p{.68in}|p{.68in}|p{.68in}}0&1&2&3&4&5&6\\\hline&&&&&&\\7&8&9&10&11&12&13\\\hline&&&&&&\\\end{tabular}\end{center}
\Task{3. Justify each mark.}{To label W, draw one legal move to an L. To label L, check \textbf{every} legal move and show it leads to a W. Start with 0 and work upward.}
\WriteBox{One W with its winning move; one L with all its replies:}{1.3in}
\StudentFooter
\newpage
\PageTitle{GRADES 2--3}{One good move, or every bad move?}\NameLine
\RuleBox{Two players alternate. Take \textbf{1, 3, or 4} counters each turn. Taking 2 is illegal. Never take more than remain. Taking the last counter wins.}\Task{4. Two claims to test.}{Alex says: ``From 8, take 3 because taking more is better.'' Bea says: ``From 10, there is only one winning move.'' Use your chart to check both.}
\WriteBox{Alex's move and a better move, with reasons:}{1in}
\WriteBox{Every winning move from 10:}{.8in}
\Task{5. Make a strategy card.}{Choose a pile between 8 and 14. If it is W, show a winning first move and how you will use your chart after that. If it is L, explain why every first move gives away the advantage.}
\WriteBox{My pile and my strategy certificate:}{1.2in}
\textbf{Go further:} Mark 14 through 20. Do you see a repeating pattern? What would you need to prove before trusting it for a pile of 100?\StudentFooter
\end{document}
